\documentclass[10pt,a4paper]{article}

\usepackage[
    left=1.35cm,
    right=1.35cm,
    top=1.15cm,
    bottom=1.15cm
]{geometry}

\usepackage{enumitem}
\usepackage{hyperref}
\usepackage{xcolor}
\usepackage{titlesec}
\usepackage{fontspec}

\definecolor{darkblue}{HTML}{0F172A}
\definecolor{accentblue}{HTML}{2563EB}

\hypersetup{
    colorlinks=true,
    urlcolor=accentblue,
    linkcolor=darkblue
}

% =========================================================
% GENERAL
% =========================================================

\setlength{\parindent}{0pt}
\setlength{\parskip}{0pt}
\renewcommand{\familydefault}{\sfdefault}
\linespread{1.02}

\raggedbottom

% =========================================================
% SECTION STYLE
% =========================================================

\titleformat{\section}
    {\normalsize\bfseries\color{darkblue}}
    {}
    {0pt}
    {}
    [\titlerule]

\titlespacing*{\section}
    {0pt}
    {0pt}
    {0pt}

% =========================================================
% LISTS
% =========================================================

\setlist[itemize]{
    leftmargin=1.15em,
    itemsep=0pt,
    topsep=0pt,
    parsep=0pt,
    partopsep=0pt
}

\setlist[description]{
    style=multiline,
    leftmargin=3.35cm,
    labelwidth=3.35cm,
    labelsep=0.12cm,
    itemsep=0pt,
    topsep=0pt,
    parsep=0pt,
    partopsep=0pt
}

\begin{document}

% =========================================================
% HEADER
% =========================================================

\begin{center}

    {\LARGE \textbf{ANDRÉ NICOLAS SILVA}} \\[-2pt]

    {\normalsize \textbf{Analista de Negócios | Comercial Tech | Dados e Processos}} \\[1pt]

    \footnotesize
    Maceió, AL \quad • \quad
    \href{mailto:devandrenicolas@gmail.com}{devandrenicolas@gmail.com} \\[-1pt]

    \textbf{\href{https://portfolio-andrenicolas.vercel.app}{Portfólio}}
    \quad • \quad
    \href{https://github.com/devAndreNicolas}{GitHub}
    \quad • \quad
    \href{https://www.linkedin.com/in/devandrenicolas}{LinkedIn}

\end{center}

\vspace{-3pt}

% =========================================================
% SUMMARY
% =========================================================

\section*{RESUMO PROFISSIONAL}

\footnotesize

Profissional de tecnologia com formação em Sistemas para Internet e experiência na análise de necessidades de produto, regras de negócio e dados em uma plataforma SaaS usada por \textbf{10+ empresas}. Participa da tradução de necessidades de usuários e áreas de produto em fluxos, requisitos e soluções técnicas, acompanhando implementação, validação e evolução em produção.

Atua com \textbf{SQL/PostgreSQL}, análise de dados de aplicações, eventos e funis de produto no PostHog, além de competências em \textbf{Excel, Microsoft Office e Power BI para análise e relatórios}. Combina contexto de negócio com conhecimento técnico para apoiar decisões de produto e colaborar com engenharia. O cargo exercido na MSPA é \textbf{Frontend Engineer}; esta variante direciona competências transferíveis para análise de negócios e Comercial Tech.
% =========================================================
% SKILLS
% =========================================================

\section*{COMPETÊNCIAS PRINCIPAIS}

\footnotesize

\begin{description}
    \item[\textbf{Análise de Negócios \& Processos}]
    Levantamento e refinamento de requisitos, regras de negócio, análise de fluxos, identificação de necessidades e definição de escopo com produto e engenharia

    \item[\textbf{Dados \& Indicadores de Produto}]
    SQL, PostgreSQL, consultas e investigação de dados; eventos, funis e feedback de usuários no PostHog

    \item[\textbf{Excel, Office \& Power BI}]
    Microsoft Office, Excel para organização e análise de dados e Power BI para análise e elaboração de relatórios

    \item[\textbf{Relatórios \& Automação}]
    Análise de informações, documentação técnica, automações entre sistemas e workflow de engenharia assistido por IA

    \item[\textbf{Colaboração e Entrega}]
    Comunicação com produto e engenharia, validação de soluções, testes, acompanhamento de mudanças e trabalho colaborativo

    \item[\textbf{Base Técnica}]
    APIs REST e internas, integrações, Angular, TypeScript, Go, Node.js, Java/Spring em contribuições na MSPA; ferramentas de IA e checks determinísticos
\end{description}

% =========================================================
% EXPERIENCE
% =========================================================

\section*{EXPERIÊNCIA PROFISSIONAL}

\noindent
\textbf{MSPA — Frontend Engineer (contribuições full stack)}
\hfill
\textit{Abr 2025 -- Presente}

\textit{Plataforma SaaS para Privacidade, Compliance (LGPD) e Auditoria Digital}

\begin{itemize}

    \item
    \textbf{Produto \& Tomada de Decisões Técnicas:}
    Traduziu requisitos de negócio e necessidades de produto em
    soluções técnicas, contribuindo para escopo de features,
    regras de negócio, UX, arquitetura, estratégia de
    implementação e trade-offs técnicos.

    \item
    \textbf{Levantamento de Necessidades e Requisitos:}
    Participou da evolução do Compass e do Trust Center, produtos
    SaaS usados por \textbf{10+ empresas}, conectando necessidades
    de usuários e produto a regras de negócio, escopo, fluxos e
    soluções implementadas em produção.

    \item
    \textbf{Análise de Produto e Processos:}
    Traduziu requisitos de negócio em soluções técnicas e
    contribuiu para avaliar fluxos, prioridades, regras,
    alternativas de implementação e trade-offs com produto e
    engenharia.

    \item
    \textbf{Análise de Dados e Comportamento:}
    Participou de entrevistas e validações de UX, instrumentação
    de eventos e análise de funis no \textbf{PostHog}, usando
    feedback e comportamento observado para apoiar decisões sobre
    os fluxos do Compass.

    \item
    \textbf{SQL e Investigação de Dados:}
    Trabalhou com \textbf{SQL e PostgreSQL}, consultas,
    persistência, migrações e investigação de problemas de dados
    relacionados às regras e aos fluxos da aplicação.

    \item
    \textbf{Soluções e Automação:}
    Contribuiu com integrações e automações de sistemas, testes e
    validação de mudanças; utilizou ferramentas de IA no workflow
    de engenharia com contexto estruturado e checks de qualidade.

    \item
    \textbf{Colaboração entre Áreas:}
    Trabalhou com produto e engenharia durante o refinamento,
    implementação e validação de soluções, mantendo alinhamento
    sobre escopo e comportamento esperado.

    \item
    \textbf{Coordenação de Workflow:}
    Reorganizou o fluxo de desenvolvimento em fase crítica de um
    projeto com equipe de \textbf{6 pessoas}, coordenando
    diretamente 4 integrantes para melhorar clareza de escopo e
    organização das entregas.

\end{itemize}

% =========================================================
% PROJECTS
% =========================================================

\section*{PROJETOS RELEVANTES}

\noindent
\textbf{RendeCerto — Simulador de cenários financeiros}
\hfill
\textit{Next.js, TypeScript, Vitest}

\footnotesize

Implementou regras determinísticas de cálculo e comparação de cenários financeiros, com testes automatizados para validar resultados em diferentes entradas.

\vspace{0.5pt}

\noindent
\textbf{Compass — Produto SaaS e análise de fluxos}
\hfill
\textit{Angular, TypeScript, Lit, Web Components, Cloudflare}

\footnotesize

Participou da análise de fluxos de produto, entrevistas e validações de UX, análise de eventos e funis no PostHog e interpretação de feedback para apoiar decisões sobre a experiência do Compass.

% EDUCATION
% =========================================================

\section*{FORMAÇÃO}

\noindent
\textbf{Tecnólogo em Sistemas para Internet} — UNCISAL
\hfill
\textit{Concluído em 2026}

% =========================================================
% LANGUAGES
% =========================================================

\vspace{-2pt}
\section*{IDIOMAS}

\noindent
\textbf{Português:} Nativo
\quad • \quad
\textbf{Inglês:} Intermediário/Avançado (em aprimoramento contínuo)

\end{document}
